\section*{Chapter \ref{ch:g6:exploring-our-environment} --- Exploring Our Environment}

\begin{solution}{exo:g6:exploring-our-environment:1}
\omterm{def:g6:exploring-our-environment:components}{Living things} (moss, woodlice); \omterm{def:g6:exploring-our-environment:components}{mineral matter} (the wall's
stone, the puddle); traces and remains of \omterm{def:g6:exploring-our-environment:components}{living things} (dead
\omterm{def:g1:plants-around-us:parts}{leaves}, the snail \omterm{def:g1:animals-around-us:covering}{shell}).
\end{solution}

\begin{solution}{exo:g6:exploring-our-environment:2}
\omterm{def:g6:exploring-our-environment:conditions}{Temperature}, with a thermometer (in degrees Celsius);
\omterm{def:g6:exploring-our-environment:conditions}{humidity} of air and soil; \omterm{def:g6:exploring-our-environment:conditions}{light}, from full sun to deep shade.
\end{solution}

\begin{solution}{exo:g6:exploring-our-environment:3}
Living: the fern, the centipede. Mineral: the puddle, the
wall's stone. Traces: the snail \omterm{def:g1:animals-around-us:covering}{shell}, the bird droppings.
\end{solution}

\begin{solution}{exo:g6:exploring-our-environment:4}
Sixteen degrees (\qty{38}{\degreeCelsius} against
\qty{22}{\degreeCelsius}) and full sun against deep shade
(dry against damp, too). The drought-and-glare weed matches
the south foot; moss, \omterm{prop:g4:plants-without-seeds:damp}{ferns} and woodlice match the cool damp
corner.
\end{solution}

\begin{solution}{exo:g6:exploring-our-environment:5}
Its covering lets water escape, so in dry sun its body loses
its water fatally fast. That is why woodlice are found in
damp, dark \omterm{def:g6:exploring-our-environment:microhabitat}{microhabitats} --- under stones, under dead wood.
\end{solution}

\begin{solution}{exo:g6:exploring-our-environment:6}
For example: a nibbled hazelnut --- the wood mouse; a
spider's web --- its spider; tracks in the mud --- the heron
(droppings, \omterm{def:g1:animals-around-us:covering}{feathers}, \omterm{def:g1:animals-around-us:covering}{shells} serve the same way).
\end{solution}

\begin{solution}{exo:g6:exploring-our-environment:7}
A closed box, half damp paper and half dry, woodlice released
along the middle line; minutes later nearly all sit on the
damp side (likewise shade against \omterm{def:g6:exploring-our-environment:conditions}{light}). It shows each
\omterm{def:g1:animals-around-us:animal}{animal} settles where \omterm{def:g6:exploring-our-environment:conditions}{conditions} suit its body --- the group's
map draws itself with no herding.
\end{solution}

\begin{solution}{exo:g6:exploring-our-environment:8}
Because only measurements made the same way can be compared:
a difference must come from the sites, not from the
measuring. It is the field version of the fair test --- one
thing varies (the site), everything else is kept equal.
\end{solution}

\begin{solution}{exo:g6:exploring-our-environment:9}
Stone top: lit, dry, warm --- lichen (with the sunning fly as
a visitor). Stone underside: dark, damp, cool --- woodlice
(or the centipede). Five centimetres apart, two matched sets
of inhabitants.
\end{solution}

\begin{solution}{exo:g6:exploring-our-environment:10}
Step 1: mark a definite square metre. Step 2: measure ---
thermometer reading written down, soil dampness felt and
noted, \omterm{def:g6:exploring-our-environment:conditions}{light} class noted. Step 3: search properly, stones
lifted and replaced, all three component kinds listed. Step
4: record on the spot --- counts, not ``some''. Step 5: no
comparison was even possible with one site; survey a second
site the same way.
\end{solution}

\begin{solution}{exo:g6:exploring-our-environment:11}
Explanation: ivy thrives in the north wall's range of
\omterm{def:g6:exploring-our-environment:conditions}{conditions} (shade, dampness) and not in the south wall's
(glare, drought). Test: measure both wall feet ---
\omterm{def:g6:exploring-our-environment:conditions}{temperature}, dampness, \omterm{def:g6:exploring-our-environment:conditions}{light} --- and check ivy's known
preferences; or find ivy on other walls and record which
orientation its walls share.
\end{solution}

\begin{solution}{exo:g6:exploring-our-environment:12}
Remove the wall difference: run the choice in a round box
(wall everywhere the same), with the damp and dry halves
swapped between runs --- damp east in run 1, damp west in
run 2. If the woodlice follow the dampness wherever it is
put, walls are ruled out; only the tested condition changed
between runs, everything else held equal.
\end{solution}

\begin{solution}{pb:g6:exploring-our-environment:1}
\textbf{1.} The tall hedge: it shades ditch N (and shelters
it), while ditch S faces the sun all day.

\textbf{2.} \omterm{def:g6:exploring-our-environment:conditions}{Temperature}: \qty{19}{\degreeCelsius} against
\qty{31}{\degreeCelsius}. \omterm{def:g6:exploring-our-environment:conditions}{Humidity}: ditch N's soil is damp
(sticks), ditch S's is dry (dust). \omterm{def:g6:exploring-our-environment:conditions}{Light} completes the trio:
shade against full sun.

\textbf{3.} \omterm{def:g6:exploring-our-environment:conditions}{Conditions} change through the day: by 6 p.m. any
site is cooler and its \omterm{def:g6:exploring-our-environment:conditions}{light} lower, so the difference
measured would mix site with hour. Same hour keeps the
comparison fair.

\textbf{4.} So that the search effort is equal: more area or
more time at one site would find more inhabitants there ---
an unfair test of which ditch is richer.

\textbf{5.} Both lists name only \omterm{def:g6:exploring-our-environment:components}{living things}. Missing:
the mineral components (soil, stones, water) and the traces
--- which parts I of a proper inventory must also record.

\textbf{6.} Woodlice: ditch N --- they need dampness and
shade or they dry out. The toad: ditch N --- moist \omterm{def:g1:the-five-senses:sense}{skin},
needing damp shelter. The lizard: ditch S --- it warms its
body in the sun, hunting where it can bask.

\textbf{7.} Either snails live there hidden (perhaps active
only in damp night hours), or the \omterm{def:g1:animals-around-us:covering}{shell} is a remain from
elsewhere or another season --- a trace proves presence
\emph{at some time}, not residence now.

\textbf{8.} Dark, damp, cool. The rule of
\cref{prop:g6:exploring-our-environment:decide}: each spot
is inhabited by the \omterm{def:g5:biodiversity:species}{species} whose ranges its \omterm{def:g6:exploring-our-environment:conditions}{conditions}
match.

\textbf{9.} N has moss and toads because its shade keeps it
cool and damp, which their bodies require; S has lizards and
waxy-leaved weeds because its full sun makes it hot and dry,
which suits a basking hunter and a drought-proof \omterm{def:g1:plants-around-us:plant}{plant}.

\textbf{10.} Without the hedge, ditch N gets full sun: hotter,
drier --- \omterm{def:g6:exploring-our-environment:conditions}{conditions} drifting toward today's ditch S. Expect
the damp-lovers (moss, \omterm{prop:g4:plants-without-seeds:damp}{ferns}, woodlice, the toad) to dwindle
and sun-lovers (dry grasses, grasshoppers, perhaps the
lizard) to move in.

\textbf{11.} Because the re-survey must differ from the first
in one thing only --- the hedge's absence. New sites,
instruments or season would add differences of their own,
and the change could no longer be pinned on the hedge: the
fair-test rule at the scale of a landscape.

\textbf{12.} Each woodlouse wanders and turns until its body
is comfortable, so the group ends gathered where dampness
is --- the \omterm{def:g1:animals-around-us:animal}{animals}' positions record the \omterm{def:g6:exploring-our-environment:conditions}{conditions}, as
surely as thermometer readings. Nobody herds them; each
individual's own comfort does the map-drawing, which is
exactly why inhabitants can be read as measurements.
\end{solution}
